\documentclass[11pt]{article}
\usepackage{amsmath,amssymb,amsthm}
\usepackage[margin=1in]{geometry}

\newtheorem{theorem}{Theorem}
\newtheorem{lemma}[theorem]{Lemma}

\title{Problem 518: Prime Triples and Geometric Sequences}
\author{Project Euler Solutions}
\date{}

\begin{document}
\maketitle

\section{Problem Statement}

Let $S(n)$ be the sum of $a + b + c$ over all triples $(a, b, c)$ of primes with $a < b < c < n$ such that $a+1, b+1, c+1$ form a geometric sequence. Given $S(100) = 1035$, find $S(10^8)$.

\section{Mathematical Foundation}

\begin{theorem}[Geometric Sequence Characterization]
Positive integers $A < B < C$ form a geometric sequence if and only if $B^2 = AC$.
\end{theorem}

\begin{proof}
If the ratio is $r > 1$, then $B = Ar$, $C = Ar^2$, so $B^2 = A^2r^2 = AC$. Conversely, $B^2 = AC$ with $A < B < C$ implies $r = B/A > 1$ and $C = Ar^2$.
\end{proof}

\begin{theorem}[Parametrization]
Every valid triple can be written as
\[
a+1 = ky^2, \quad b+1 = kxy, \quad c+1 = kx^2,
\]
where $k \geq 1$, $x > y \geq 1$, and $\gcd(x,y) = 1$.
\end{theorem}

\begin{proof}
Write the common ratio $r = (b+1)/(a+1) = x/y$ in lowest terms. Then $a+1 = ky^2$ for some integer $k$, and the expressions for $b+1$ and $c+1$ follow by multiplying by $x/y$ and $(x/y)^2$ respectively. Since $\gcd(x,y) = 1$, all values are integers.
\end{proof}

\begin{lemma}[Parameter Bounds]
For $c < n$, we need $kx^2 \leq n$, giving $x \leq \sqrt{n}$ and $k \leq n/x^2$.
\end{lemma}

\begin{proof}
From $c + 1 = kx^2 \leq n$, the bounds follow immediately.
\end{proof}

\section{Algorithm}

\begin{enumerate}
    \item Sieve all primes up to $n = 10^8$ using the Sieve of Eratosthenes.
    \item For each $x$ from 2 to $\lfloor\sqrt{n}\rfloor$:
    \begin{enumerate}
        \item For each $y$ from 1 to $x - 1$ with $\gcd(x, y) = 1$:
        \begin{enumerate}
            \item For each $k$ with $kx^2 \leq n$:
            \begin{enumerate}
                \item Compute $a = ky^2 - 1$, $b = kxy - 1$, $c = kx^2 - 1$.
                \item If all three are prime and $a \geq 2$, add $a + b + c$ to the sum.
            \end{enumerate}
        \end{enumerate}
    \end{enumerate}
\end{enumerate}

\section{Complexity Analysis}

\begin{itemize}
    \item \textbf{Time:} $O(n \log\log n)$ for the sieve. The triple enumeration is $O(n \log n)$ in the worst case but much less in practice due to $\gcd$ and primality filtering.
    \item \textbf{Space:} $O(n)$ for the prime sieve.
\end{itemize}

\section{Answer}

\[
\boxed{100315739184392}
\]

\end{document}
